\documentclass[12pt, leqno]{article}
\usepackage[english,russian]{babel}
\usepackage{fontspec}
\usepackage[utf8]{inputenc}
\usepackage{epigraph}
\usepackage{tabularx}
\usepackage[leqno]{amsmath}
\setmainfont{Times New Roman}
\setlength\parindent{0pt}
\usepackage[colorlinks=true,urlcolor=blue]{hyperref}        
\usepackage{minted} 
\usepackage[xetex,final]{graphicx}
\usepackage{float}%"Плавающие" картинки
\usepackage{wrapfig}%Обтекание фигур (таблиц, картинок и прочег
\usepackage{epstopdf}
\usepackage{amssymb}
 
\begin{document}

\section*{Биномиальные коээфициэнты. Рекурсия}
\begin{flushleft}
\textbf{Вопросы рекурсии:} Рекурентные формулы и реализация
\end{flushleft}


Биномиальные коэффициенты часто обозначаются   $ {\binom {n}{k}} $ 
и читаются и читаются как «число сочетаний из $ n $ элементов по $ k $


вот эта формула в коде взята отсюда:
\begin{flushleft}
{\binom {n}{k}} = \binom {n}{k - 1}\times {\frac {n+1-k}{k}}
\end{flushleft}

\clearpage

\noindent Образец кода на haskell :  
\begin{minted}[fontsize=\small]{haskell}

module Math where

{- Binomial cooeficient Newtone / pascal  triangle -}
{- http://mech.math.msu.su/~shvetz/54/inf/perl-examples/PerlExamples_PascalTriangle.xhtml -}
{- https://stackoverflow.com/questions/42798257/add-a-element-at-the-end-of-list-in-haskell -}
{- https://rosettacode.org/wiki/List_comprehensions -}

pascal :: Integral a => a -> a -> a  
pascal 0 0 = 1
pascal n 0 = 1
pascal n r | n == r = 1
           | otherwise  = pascal (n - 1) (r - 1) + pascal (n - 1) r

binom :: Integral a => a -> [a]
binom n = map (pascal n) [0..n]

makeTriangle :: Integer -> [[Integer]]
makeTriangle  n = [ (binom x) | x <- [1..n]]


module Main where

import Math

{- https://libeldoc.bsuir.by/bitstream/123456789/56340/1/Mychko_Treugol%27nik.pdf -}
{- http://mech.math.msu.su/~shvetz/54/inf/perl-examples/PerlExamples_PascalTriangle.xhtml -}
{- http://mech.math.msu.su/~vvb/Dush/Haskell/index.html -}
{- https://gist.github.com/gubatron/c690b46f72e0bb1bab378da8a49348c1 -}

listToString = unwords. map show
printTriangle n = mapM_ putStrLn (map listToString (makeTriangle 10))

main = printTriangle 10



\end{minted}  

результат нужно получить самостоятельно:
    
\end{document}